% Source: Weinberg GR, physical PDF pp. 217--224
%         (printed pp. 194--201).
% Coverage: Section 8.6, Bound Orbits: Precession of Perihelia; equations
%           (8.6.1)--(8.6.13), Figure 8.2, Table 8.3.
% Figures/tables/footnotes: Figure 8.2; Table 8.3; reference markers 5--12.
% Status: source-reviewed and compile-clean.
% Uncertainties: SOURCE ERRATUM -- the scan says "Table 8.2" where the data
%                are in Table 8.3; the reference is corrected below.

\subsection{Bound Orbits: Precession of Perihelia}\label{sec:8.6}

Now consider a test particle bound in an orbit around the sun. (See
Figure~\ref{fig:8.2}.) At perihelia and aphelia, \(r\) reaches its minimum
and maximum values \(r_-\) and \(r_+\), and at both points
\(\dd r/\dd\varphi\) vanishes, so Eq.~\eqref{eq:8.4.29} gives
\[
  \frac{1}{r_+^2}-\frac{1}{J^2B(r_+)}=-\frac{E}{J^2},
  \qquad
  \frac{1}{r_-^2}-\frac{1}{J^2B(r_-)}=-\frac{E}{J^2}.
\]

% Source: Weinberg GR, physical PDF p. 218 (printed p. 195).
% Coverage: Figure 8.2, geometric elements of an eccentric ellipse.
% Status: vector reconstruction; source-reviewed and compile-clean.
% Uncertainties: none.

\begingroup
\renewcommand{\thefigure}{8.2}
\begin{figure}[t]
\centering
\begin{tikzpicture}[x=0.82cm,y=0.82cm,>=Latex]
  \draw[line width=0.75pt] (0,0) ellipse (4.4 and 2.25);
  \coordinate (F) at (3.35,0);
  \fill (F) circle (1.8pt);

  % Aphelion and perihelion distances measured from the occupied focus.
  \draw[<->] (-4.4,0) -- (F);
  \node[above] at (-0.55,0) {$r_+$};
  \draw[<->] (F) -- (4.4,0);
  \node[below] at (3.88,-0.02) {$r_-$};

  % Semilatus rectum at the focus.
  \draw[<->] (F) -- (3.35,1.45);
  \node[left] at (3.28,0.83) {$L$};

  % Semimajor axis and focal offset above the ellipse.
  \draw (0,2.45) -- (0,4.18);
  \draw (4.4,2.28) -- (4.4,4.18);
  \draw[<->] (0,3.92) -- (4.4,3.92);
  \node[above] at (2.2,3.92) {$a$};

  \draw (3.35,2.10) -- (3.35,3.35);
  \draw[<->] (0,3.10) -- (3.35,3.10);
  \node[above] at (1.68,3.10) {$ea$};

  % Semiminor axis.
  \draw (4.95,0) -- (5.95,0);
  \draw (0,2.25) -- (5.95,2.25);
  \draw[<->] (5.65,0) -- (5.65,2.25);
  \node[right] at (5.75,1.15) {$a\sqrt{1-e^2}$};
\end{tikzpicture}
\caption{Elements of an ellipse referred to in the calculation of the
precession of planetary orbits. (The ellipse here has the same eccentricity
as the orbit of Icarus.)}
\label{fig:8.2}
\end{figure}
\endgroup


From these two equations we can derive values for the two constants of the
motion:
\begin{equation}
  E
  =
  \frac{r_+^2/B(r_+)-r_-^2/B(r_-)}
       {r_+^2-r_-^2},
  \tag{8.6.1}\label{eq:8.6.1}
\end{equation}
\begin{equation}
  J^2
  =
  \frac{B^{-1}(r_+)-B^{-1}(r_-)}
       {r_+^{-2}-r_-^{-2}}.
  \tag{8.6.2}\label{eq:8.6.2}
\end{equation}
The angle swept out by the position vector as \(r\) increases from \(r_-\)
is given by Eq.~\eqref{eq:8.4.30} as
\[
  \varphi(r)
  =
  \varphi(r_-)
  +\int_{r_-}^{r}
  A^{1/2}(r)
  \left[
    \frac{1}{J^2B(r)}-\frac{E}{J^2}-\frac{1}{r^2}
  \right]^{-1/2}\frac{\dd r}{r^2},
\]
or, using Eqs.~\eqref{eq:8.6.1} and \eqref{eq:8.6.2},
\begin{equation}
  \begin{split}
  \varphi(r)-\varphi(r_-)
  =\int_{r_-}^{r}
  \biggl[
  &\frac{
    r_-^2\!\left[B^{-1}(r)-B^{-1}(r_-)\right]
    -r_+^2\!\left[B^{-1}(r)-B^{-1}(r_+)\right]
  }{
    r_+^2r_-^2\!\left[B^{-1}(r_+)-B^{-1}(r_-)\right]
  }
  -\frac{1}{r^2}
  \biggr]^{-1/2}\\
  &{}\times A^{1/2}(r)r^{-2}\,\dd r.
  \end{split}
  \tag{8.6.3}\label{eq:8.6.3}
\end{equation}
The change in \(\varphi\) as \(r\) decreases from \(r_+\) to \(r_-\) is
the same as the change in \(\varphi\) as \(r\) increases from \(r_-\) to
\(r_+\), so the total change in \(\varphi\) per revolution is
\(2|\varphi(r_+)-\varphi(r_-)|\). This would equal \(2\pi\) if the orbit
was a closed ellipse, so in general the orbit precesses in each revolution
by an angle
\begin{equation}
  \Delta\varphi
  =
  2|\varphi(r_+)-\varphi(r_-)|-2\pi.
  \tag{8.6.4}\label{eq:8.6.4}
\end{equation}

Using the exact values of \(A(r)\) and \(B(r)\) given by the Schwarzschild
solution \eqref{eq:8.2.10}, \eqref{eq:8.2.11} in
Eq.~\eqref{eq:8.6.3} would yield formulas for \(\varphi(r)\) and
\(\Delta\varphi\) as elliptic integrals, and to evaluate them numerically
we would have to expand in \(MG/r_-\) and \(MG/r_+\). Instead we shall
expand in the integrand, using for \(A(r)\) and \(B(r)\) the Robertson
expansions \eqref{eq:8.3.8}, \eqref{eq:8.3.9}:
\begin{equation}
  \begin{aligned}
  A(r)&=1+2\gamma\frac{MG}{r}+\cdots,\\
  B(r)&=1-\frac{2MG}{r}
  +2(\beta-\gamma)\frac{M^2G^2}{r^2}+\cdots.
  \end{aligned}
  \tag{8.6.5}\label{eq:8.6.5}
\end{equation}
Note that there is a complete cancellation in Eq.~\eqref{eq:8.6.3} of
the leading term in \(B(r)\) but not of that in \(A(r)\), so to calculate
\(\varphi\) and \(\Delta\varphi\) to first order in \(MG/r_\pm\) we need
\(B(r)\) to second order in \(MG/r\), whereas \(A(r)\) will be needed only
to first order.

It saves a great deal of work if we realize that by using the expansion
\[
  B^{-1}(r)
  \simeq
  1+\frac{2MG}{r}
  +2(2-\beta+\gamma)\frac{M^2G^2}{r^2},
\]
we make the argument of the first square root in
Eq.~\eqref{eq:8.6.3} a quadratic function of \(1/r\). Furthermore, it
vanishes at \(r=r_\pm\), so
\begin{equation}
  \begin{split}
  &\frac{
    r_-^2\!\left[B^{-1}(r)-B^{-1}(r_-)\right]
    -r_+^2\!\left[B^{-1}(r)-B^{-1}(r_+)\right]
  }{
    r_+^2r_-^2\!\left[B^{-1}(r_+)-B^{-1}(r_-)\right]
  }
  -\frac{1}{r^2}\\
  &\hspace{4cm}
  =C\left(\frac1{r_-}-\frac1r\right)
  \left(\frac1r-\frac1{r_+}\right).
  \end{split}
  \tag{8.6.6}\label{eq:8.6.6}
\end{equation}
The constant \(C\) can be determined by letting \(r\to\infty\):
\[
  C
  =
  \frac{
    r_+^2[1-B^{-1}(r_+)]
    -r_-^2[1-B^{-1}(r_-)]
  }{
    r_+r_-[B^{-1}(r_+)-B^{-1}(r_-)]
  },
\]
or factoring out of numerator and denominator a common factor
\(2(r_--r_+)MG\),
\begin{equation}
  C
  \simeq
  1-(2-\beta+\gamma)MG
  \left(\frac1{r_+}+\frac1{r_-}\right).
  \tag{8.6.7}\label{eq:8.6.7}
\end{equation}
Using Eqs.~\eqref{eq:8.6.5}--\eqref{eq:8.6.7} in
Eq.~\eqref{eq:8.6.3} gives then
\[
  \begin{split}
  \varphi(r)-\varphi(r_-)
  \simeq{}&
  \left[
    1+\frac12(2-\beta+\gamma)MG
    \left(\frac1{r_+}+\frac1{r_-}\right)
  \right]\\
  &{}\times
  \int_{r_-}^{r}
  \frac{[1+\gamma MG/r]\,\dd r}
  {r^2[(1/r_--1/r)(1/r-1/r_+)]^{1/2}}.
  \end{split}
\]
The integral is made trivial by introducing a new variable \(\psi\):
\begin{equation}
  \frac1r
  \equiv
  \frac12\left(\frac1{r_+}+\frac1{r_-}\right)
  +\frac12\left(\frac1{r_+}-\frac1{r_-}\right)\sin\psi.
  \tag{8.6.8}\label{eq:8.6.8}
\end{equation}
We then find
\begin{equation}
  \begin{split}
  \varphi(r)-\varphi(r_-)
  ={}&
  \left[
    1+\frac12(2-\beta+2\gamma)MG
    \left(\frac1{r_+}+\frac1{r_-}\right)
  \right]
  \left(\psi+\frac{\pi}{2}\right)\\
  &-\frac12\gamma MG
  \left(\frac1{r_+}-\frac1{r_-}\right)\cos\psi.
  \end{split}
  \tag{8.6.9}\label{eq:8.6.9}
\end{equation}
At aphelion \(\psi=\pi/2\), so Eqs.~\eqref{eq:8.6.4} and
\eqref{eq:8.6.9} give the precession per revolution as
\begin{equation}
  \Delta\varphi
  =
  \left(\frac{6\pi MG}{L}\right)
  \left(\frac{2-\beta+2\gamma}{3}\right)
  \qquad\text{(radians/revolution)},
  \tag{8.6.10}\label{eq:8.6.10}
\end{equation}
where \(L\) is the dimension of the ellipse called the \emph{semilatus
rectum},
\[
  \frac1L
  \equiv
  \frac12\left(\frac1{r_+}+\frac1{r_-}\right).
\]
The elements of planetary orbits usually found in tables are the semimajor
axis \(a\) and eccentricity \(e\), defined by
\[
  r_\pm=(1\pm e)a.
\]
Hence we can determine \(L\) from \(a\) and \(e\) by using the formula
\[
  L=(1-e^2)a.
\]
Einstein's field equations yield \(\beta=\gamma=1\), so they predict a
precession
\begin{equation}
  \Delta\varphi
  =
  6\pi\frac{MG}{L}
  \qquad\text{radians/revolution}.
  \tag{8.6.11}\label{eq:8.6.11}
\end{equation}
This is positive, meaning that the whole orbit should precess in the same
direction as the motion of the test particle. In the Brans--Dicke theory
Eqs.~\eqref{eq:8.6.10} and \eqref{eq:8.3.3} give
\begin{equation}
  \Delta\varphi
  =
  \left(\frac{6\pi MG}{L}\right)
  \left(\frac{3\omega+4}{3\omega+6}\right).
  \tag{8.6.12}\label{eq:8.6.12}
\end{equation}

Once again, we should ask ourselves what the predicted value of
\(\Delta\varphi\) means. This is not a scattering experiment like the
deflection of light by the sun; here we are dealing with an object that
never gets out to infinity where the metric is Minkowskian. Any observation
of the test particle's motion by optical or radar astronomers will make use
of light rays that are themselves affected by the gravitational field, and
if careful corrections are not made for the deflection of light, the
astronomer's reported \(\varphi(r)\) values will, at any given radius \(r\),
contain errors of the order of \(MG/L\). [See Eq.~\eqref{eq:8.5.8}.]
However, in practice these fine points do not really matter, because the
precession is \emph{cumulative}. Equation~\eqref{eq:8.6.10} shows that
after \(N\) revolutions the perihelia will have advanced by an angle of
order \(NMG/L\), so if \(N\gg1\) it is unnecessary to worry about an error
in \(\varphi\) of order \(MG/L\). Indeed, Eq.~\eqref{eq:8.6.11} tells us
that the perihelion will return to its original azimuth after
\(L/3MG\gg1\) revolutions, a prediction that clearly has nothing to do with
how we define \(r\) or \(\varphi\).

For Mercury we must take \(L=55.3\times10^6\,\mathrm{km}\), and of course
\(MG=1.475\,\mathrm{km}\), so Eq.~\eqref{eq:8.6.11} gives
\(\Delta\varphi=0.1038''\) per revolution. Since Mercury makes 415
revolutions per century, the prediction of general relativity is that
\[
  \Delta\varphi=43.03''\text{ per century}\quad(\text{Mercury}).
\]
Fortunately there are accurate observations of Mercury going back to 1765.
These data were reanalyzed by
Clemence\hyperref[ch8-ref-5]{\textsuperscript{5}} in 1943; he finds
\(\Delta\varphi=43.11\pm0.45''\) per century, essentially confirming
Newcomb's earlier value (see Section~\ref{sec:1.2}) and in excellent
agreement with general relativity. Taken at face value this agreement shows
that the correction factor in Eq.~\eqref{eq:8.6.11} is
\[
  \left(\frac{2-\beta+2\gamma}{3}\right)=1.00\pm0.01.
\]
This is by far the most important experimental verification of general
relativity, both by virtue of its high accuracy, and because it alone is
sensitive to the coefficient \(\beta\) appearing in the second-order term
in \(g_{tt}\).

\begingroup
\renewcommand{\thetable}{8.3}
\begin{table}[tbp]
\centering
\caption{Comparison of Theoretical and Observed Centennial Precessions of
Planetary Orbits.\hyperref[ch8-ref-6]{\textsuperscript{6}}}
\label{tab:8.3}
\small
\begin{tabular}{@{}lrrrrrr@{}}
\toprule
Planet & \(a\) (\(10^6\) km) & \(e\) & \(6\pi MG/L\) &
Revolutions/century &
\multicolumn{2}{c}{\(\Delta\varphi\) (seconds/century)}\\
\cmidrule(l){6-7}
 & & & & & Gen. Rel. & Observed\\
\midrule
Mercury & 57.91 & 0.2056 & \(0.1038''\) & 415 & 43.03 &
  \(43.11\pm0.45\)\\
Venus & 108.21 & 0.0068 & \(0.058''\) & 140 & 8.6 &
  \(8.4\pm4.8\)\\
Earth & 149.60 & 0.0167 & \(0.038''\) & 100 & 3.8 &
  \(5.0\pm1.2\)\\
Icarus & 161.0 & 0.827 & \(0.115''\) & 89 & 10.3 &
  \(9.8\pm0.8\)\\
\bottomrule
\end{tabular}
\end{table}
\endgroup

Results\hyperref[ch8-ref-6]{\textsuperscript{6}} for Venus, the earth, and
Icarus are listed together with those for Mercury in
Table~\ref{tab:8.3}. Evidently the accuracy available from the major
planets degrades rapidly as we move away from the sun, both because the
smaller eccentricities make observation of the perihelia more uncertain,
and because as \(L\) increases, the precession per revolution and the
revolutions per century both decrease. Icarus was only discovered in 1949,
but is in some respects the most useful object to study because its small
size, its close approach to the earth, and the large eccentricity of its
orbit allow its precession to be determined with high accuracy. Suggestions
have been made to put an artificial satellite into an eccentric orbit close
to the sun; for instance, a satellite with \(L=10R_\odot\) would have a
centennial precession equal to \(8250''\)! The trouble here is that a small
object would be subject to nongravitational perturbations such as radiation
pressure, solar wind, and micrometeorites, which of course have a negligible
effect on Mercury and Icarus.

There are two caveats that should be kept in mind in assessing the agreement
of the observed advance of perihelia with the prediction of general
relativity. First, there are many known perturbations that contribute to the
precession of planetary orbits. In particular, Newtonian theory would give
Mercury a precession
\[
  \Delta\varphi_N=5557.62\pm0.20''\quad(\text{Mercury}),
\]
of which about \(5025''\) is due to the rotation of the earth-based
astronomical coordinate system, and about \(532''\) is due to gravitational
perturbations calculated by Newtonian perturbation theory from the motion
of the other planets, chiefly Venus, earth, and Jupiter. The precession
actually observed is
\[
  \Delta\varphi_{\mathrm{OBS}}
  =5600.73\pm0.41''\quad(\text{Mercury}),
\]
and the value \(\Delta\varphi=43.11\pm0.45''\) quoted above for the
``observed'' anomalous precession is obtained by subtracting the Newtonian
precession from what is observed, that is,
\begin{equation}
  \Delta\varphi
  =
  \Delta\varphi_{\mathrm{OBS}}-\Delta\varphi_N.
  \tag{8.6.13}\label{eq:8.6.13}
\end{equation}
One may ask how we know that this is the right quantity to compare with the
general-relativistic result of \(43.03''\) per century. That is, how do we
know that the total precession is correctly given by adding the Newtonian
value \(\Delta\varphi_N\), calculated while forgetting all effects of
general relativity, to the Einstein value \(\Delta\varphi_{\mathrm{GR}}\),
calculated forgetting all effects of planetary perturbations? To some
extent this question can be answered by noting that the general-relativistic
corrections to \(\Delta\varphi_N\) would be of order \(MG/L\) times
\(\Delta\varphi_N\), or only about \(10^{-4}\,\text{arcsec}\) per century.
For a fuller
answer we shall have to wait until the discussion in the next chapter of the
post-Newtonian approximation. But even granting that
Eq.~\eqref{eq:8.6.13} is in principle correct, we should be aware that a
very small systematic error in either \(\Delta\varphi_N\) or
\(\Delta\varphi_{\mathrm{OBS}}\) could completely destroy the agreement
between theory and observation.

The second warning is that very small \emph{unknown} effects may possibly be
contributing to the observed precession of perihelia an amount comparable
with that expected from general relativity. Indeed, we saw in Chapter~1
that Newcomb had by 1911 abandoned his earlier suggestion of a small
departure from the inverse-square law, because the observed anomalous
precession of \(43''\) per century could also be explained within Newtonian
mechanics as due to the gravitational field produced by the matter which
causes the ``zodiacal light.'' (Today we know that there is not enough
matter between Mercury and the sun to produce any appreciable precession.)
It is also possible that the sun is slightly
oblate,\hyperref[ch8-ref-7]{\textsuperscript{7}} in which case its
Newtonian potential would have an \(r^{-3}\) term, giving the planets an
anomalous precession per revolution decreasing as the inverse square of
their distance from the sun.
% SOURCE ERRATUM: the scan says Table 8.2; the data are in Table 8.3.
Table 8.3 shows that in fact the observed anomalous precession per
revolution decreases roughly as \(1/r\), not \(1/r^2\), in agreement with
the prediction of general relativity. Even more important, a large solar
oblateness would produce an anomalous precession of the planes of the inner
planets' orbits that is not
observed.\hyperref[ch8-ref-8]{\textsuperscript{8}} These two arguments
together rule out the possibility of explaining \emph{all} of the observed
anomalous precession as owing to solar oblateness, but this explanation
might account for up to 20\% of the observed effect.

To test this hypothesis Dicke and
Goldenberg\hyperref[ch8-ref-9]{\textsuperscript{9}} scanned the solar disk
photoelectrically during the period June~1 to September~23, 1966. They
concluded that the sun's polar diameter is shorter than its equatorial
diameter by \(5.0\pm0.7\) parts in \(10^5\). Taken at face value, this
would give the perihelia of Mercury an extra precession of \(3.4''\) per
century, so that only \(39.6''\) per century would be left as the
relativistic effect, an 8\% disagreement with Einstein's prediction of
\(43.03''\) per century. The Brans--Dicke theory can account for an excess
centennial precession of \(39.6''\) if we take \(\omega=6.4\). However,
there are several reasons for hesitating before we give up general
relativity.

\begin{enumerate}
  \item[(A)] To account for the solar oblateness the solar interior would
  have to be rotating about once in one or two days, very much faster than
  the observed rotation rate (i.e., once in 25 days) of the sun's surface.
  This difference in rotation rates could perhaps be
  explained\hyperref[ch8-ref-10]{\textsuperscript{10}} as due to a
  magnetic torque induced by the solar wind, which retards the rotation of
  the surface, but it is not clear that this configuration would be
  dynamically stable.\hyperref[ch8-ref-11]{\textsuperscript{11}}

  \item[(B)] Two very elaborate series of
  measurements\hyperref[ch8-ref-12]{\textsuperscript{12}} made with the
  G\"ottingen heliometer during the period 1891--1902 gave for the
  difference between the sun's equator and polar diameters the values
  \(0.36\pm0.78\) parts in \(10^5\) and \(-0.10\pm0.47\) parts in
  \(10^5\), respectively, in agreement with each other and with perfect
  sphericity, but in disagreement with the result \(+5.0\pm0.7\) parts in
  \(10^5\) of Dicke and Goldenberg. The G\"ottingen results were also
  supported by subsequent heliometer measurements. To quote
  Ashbrook:\hyperref[ch8-ref-12]{\textsuperscript{12}}
  \begin{quote}
  ``What are we to make of all this? In view of the astronomical evidence,
  can the sun's polar diameter be 0.1 seconds shorter than its equatorial,
  as Drs. Dicke and Goldenberg think? Was there some unsuspected subtle
  systematic error in the Princeton experiment? Or was there some
  unrecognized effect in all the heliometer series?''
  \end{quote}

  To do Dicke and Goldenberg justice, it should be noted that the axis of
  the oblate ellipsoid they observe follows the axis of the sun's rotation
  in its apparent annual oscillation, which certainly suggests that they
  are seeing something real.

  \item[(C)] Even if the visible solar surface is oblate, what does this
  really tell us about the shape of its mass distribution and about the
  sun's gravitational field? Dicke\hyperref[ch8-ref-7]{\textsuperscript{7}}
  argues that the observed solar surface coincides with the gravitational
  equipotential surface, but this conclusion rests on a good deal of
  astrophysical theory and could be wrong.

  \item[(D)] Finally, if Dicke and Goldenberg are right, then the 1\%
  agreement between Einstein's prediction and the observed anomalous
  precession is a mere coincidence.
\end{enumerate}
